%% Version 6.1, 1 September 2021
%
%%%%%%%%%%%%%%%%%%%%%%%%%%%%%%%%%%%%%%%%%%%%%%%%%%%%%%%%%%%%%%%%%%%%%%
% amspaperV6.tex --  LaTeX-based instructional template paper for submissions to the 
% American Meteorological Society
%
%%%%%%%%%%%%%%%%%%%%%%%%%%%%%%%%%%%%%%%%%%%%%%%%%%%%%%%%%%%%%%%%%%%%%
% PREAMBLE
%%%%%%%%%%%%%%%%%%%%%%%%%%%%%%%%%%%%%%%%%%%%%%%%%%%%%%%%%%%%%%%%%%%%%
\documentclass{ametsocV6.1}

\title{Seasonal cycle of amplified hot-day warming}

\authors{Osamu Miyawaki,\thanks{Osamu Miyawaki's current affiliation: NCAR, Boulder, Colorado},\aff{a}\correspondingauthor{Osamu Miyawaki, miyawaki@ucar.edu} 
Isla Simpson,\aff{a} 
Brian P. Medeiros,\aff{a} 
Qinqin Kong,\aff{b} 
Karen McKinnon,\aff{c} 
}

\affiliation{\aff{a}{NSF National Center for Atmospheric Research}\\
\aff{b}{Purdue University}\\
\aff{c}{University of California Los Angeles}\\
}

%%%%%%%%%%%%%%%%%%%%%%%%%%%%%%%%%%%%%%%%%%%%%%%%%%%%%%%%%%%%%%%%%%%%%
% ABSTRACT
%%%%%%%%%%%%%%%%%%%%%%%%%%%%%%%%%%%%%%%%%%%%%%%%%%%%%%%%%%%%%%%%%%%%%

\abstract{In response to anthropogenic emissions, hot days are projected to warm more than the mean day in the tropics. Climate models and theories have attributed amplified hot day warming to reduced surface specific humidity and surface evaporative cooling during those days. However, we currently do not have a complete understanding of the seasonal and spatial structure of amplified hot day warming. Here we show the spatiotemporal variation in amplified hot day warming over land can be understood based on the climatological evapotranspiration regime of hot and mean days. Specifically, hot days warm more than the mean because hot days are in the water-limited regime (evaporation is sensitive to drying soil) whereas mean days are in the energy-limited regime (evaporation is insensitive to drying soil). The results provide a physical link between amplified hot day warming and the present-day evapotranspiration regime, suggesting the possibility of an emergent constraint on amplified hot day warming.}

\begin{document}


%% Necessary!
\maketitle

%%%%%%%%%%%%%%%%%%%%%%%%%%%%%%%%%%%%%%%%%%%%%%%%%%%%%%%%%%%%%%%%%%%%%
% SIGNIFICANCE STATEMENT/CAPSULE SUMMARY
%%%%%%%%%%%%%%%%%%%%%%%%%%%%%%%%%%%%%%%%%%%%%%%%%%%%%%%%%%%%%%%%%%%%%
\statement
	 Enter significance statement here, no more than 120 words. See \url{www.ametsoc.org/index.cfm/ams/publications/author-information/significance-statements/} for details.


%%%%%%%%%%%%%%%%%%%%%%%%%%%%%%%%%%%%%%%%%%%%%%%%%%%%%%%%%%%%%%%%%%%%%
% MAIN BODY OF PAPER
%%%%%%%%%%%%%%%%%%%%%%%%%%%%%%%%%%%%%%%%%%%%%%%%%%%%%%%%%%%%%%%%%%%%%
%
\section{Introduction}

\section{Methods}

\begin{table}[t]
  \caption{List of the 15 models that comprise the CMIP6 multimodel mean of the historical and SSP3-7.0 runs.}
\begin{center}
  \renewcommand{\arraystretch}{1.0}
  \begin{tabular}{ l l }
    Model           & Ensemble run \\
    \hline
    ACCESS-CM2      & r1i1p1f1 \\
    CESM2           & r11i1p1f1 \\
    CanESM5         & r1i1p1f1 \\
    EC-Earth3       & r1i1p1f1 \\
    IPSL-CM6A-LR    & r1i1p1f1 \\
    KACE-1-0-G      & r1i1p1f1 \\
    MIROC-ES2L      & r1i1p1f2 \\
    MIROC6          & r1i1p1f1 \\
    MPI-ESM1-2-HR   & r1i1p1f1 \\
    MPI-ESM1-2-LR   & r1i1p1f1 \\
    MRI-ESM2-0      & r1i1p1f1 \\
    NorESM2-LM      & r1i1p1f1 \\
    NorESM2-MM      & r1i1p1f1 \\
    TaiESM1         & r1i1p1f1 \\
    UKESM1-0-LL     & r9i1p1f2 \\
  \end{tabular}
\end{center}
\end{table}

\section{Results}

\section{Conclusion}

\clearpage
%%%%%%%%%%%%%%%%%%%%%%%%%%%%%%%%%%%%%%%%%%%%%%%%%%%%%%%%%%%%%%%%%%%%%
% ACKNOWLEDGMENTS
%%%%%%%%%%%%%%%%%%%%%%%%%%%%%%%%%%%%%%%%%%%%%%%%%%%%%%%%%%%%%%%%%%%%%
\acknowledgments
Keep acknowledgments (note correct spelling: no ``e'' between the ``g'' and
``m'') as brief as possible. In general, acknowledge only direct help in
writing or research. Financial support (e.g., grant numbers) for the work done, 
for an author, or for the laboratory where the work was performed must be 
acknowledged here rather than as footnotes to the title or to an author's name.
Contribution numbers (if the work has been published by the author's institution 
or organization) should be placed in the acknowledgments rather than as 
footnotes to the title or to an author's name.

%%%%%%%%%%%%%%%%%%%%%%%%%%%%%%%%%%%%%%%%%%%%%%%%%%%%%%%%%%%%%%%%%%%%%
% DATA AVAILABILITY STATEMENT
%%%%%%%%%%%%%%%%%%%%%%%%%%%%%%%%%%%%%%%%%%%%%%%%%%%%%%%%%%%%%%%%%%%%%
\datastatement
The data availability statement is where authors should describe how the data underlying 
the findings within the article can be accessed and reused. Authors should attempt to 
provide unrestricted access to all data and materials underlying reported findings. 
If data access is restricted, authors must mention this in the statement.
See http://www.ametsoc.org/PubsDataPolicy for more details.


%%%%%%%%%%%%%%%%%%%%%%%%%%%%%%%%%%%%%%%%%%%%%%%%%%%%%%%%%%%%%%%%%%%%%
% APPENDIXES
%%%%%%%%%%%%%%%%%%%%%%%%%%%%%%%%%%%%%%%%%%%%%%%%%%%%%%%%%%%%%%%%%%%%%

\appendix[A] 

\appendixtitle{Title of Appendix}

\subsection*{Appendix section}




%%%%%%%%%%%%%%%%%%%%%%%%%%%%%%%%%%%%%%%%%%%%%%%%%%%%%%%%%%%%%%%%%%%%%
% REFERENCES
%%%%%%%%%%%%%%%%%%%%%%%%%%%%%%%%%%%%%%%%%%%%%%%%%%%%%%%%%%%%%%%%%%%%%
\bibliographystyle{ametsocV6}
\bibliography{references}


\end{document}
%%%%%%%%%%%%%%%%%%%%%%%%%%%%%%%%%%%%%%%%%%%%%%%%%%%%%%%%%%%%%%%%%%%%%
% END OF AMSPAPERV6.1.TEX
%%%%%%%%%%%%%%%%%%%%%%%%%%%%%%%%%%%%%%%%%%%%%%%%%%%%%%%%%%%%%%%%%%%%%
